% SUMMARY AND AUDIT TRAIL (campaign 3, 2026-08-01)
% Standalone subsection, drop-in after the three targets of
% FABLE_SECTION_orbit_energy.tex.  No new packages.
% Proved here, complete proofs: Casoratian gap identity (machine-checked
% 1500/1500 across p in {499,997,1999,4001,7919}); centered-count inequality;
% the conditional energy theorem [COINC_beta] => F_p <= (1+o(1)) p^{2-beta};
% the Parseval bridge showing COINC is the L^2(t)-average of GAP-2D-SQRT,
% hence the WEAKEST member of the sufficient-target family.
% Open and labelled conjectural: [COINC_beta] itself for every beta > 1/2.
% Numerical status points to the campaign dossier; no empirical assertion
% is used in any proof.

\subsection{A fourth target: coincidence second moments of gap Casoratians}
\label{subsec:coinc-target}

The three targets above either control every additive character twist
(GAP--2D--SQRT), the full variance of the family $\{C_h\}$ (GPRV), or the
triple returns (ATR).  This subsection isolates a fourth target, strictly
weaker than GAP--2D--SQRT at the same scale, carrying a scale dial
$\beta$: any value $\beta>\tfrac12$ already lowers the unconditional
exponent $3/2$ of Theorem~\ref{thm:oe-energy-three-halves}, and $\beta=1$
recovers E1 with an explicit constant.

For $1\le h\le M-1$ and $1\le r\le M-h$ put
\[
 \Delta_{r,h}=b_rc_{r+h}-b_{r+h}c_r\in\mathbb F_p ,
\]
the gap Casoratian of the pair of solutions at distance $h$.

\begin{lemma}[Casoratian gap identity]\label{lem:coinc-gap-identity}
For every $1\le h\le M-1$ and $1\le r\le M-h$,
\begin{equation}\label{eq:coinc-gap-identity}
 \boxed{\;\Delta_{r,h}
 =\frac{N_h(r)}{\prod_{j=1}^{h}(r+j)^3}\;}
\end{equation}
in $\mathbb F_p$; every factor $r+j$ in the range $1\le j\le h$ satisfies
$2\le r+j\le M<p$ and is invertible.  In particular, on the nonwrapping
window, $\Delta_{r,h}=0$ if and only if $\pi(r)=\pi(r+h)$, and
$\#\{r:\Delta_{r,h}=0\}=C_h$.
\end{lemma}

\begin{proof}
Fix $r$ and set $y_n=b_rc_n-c_rb_n$, a solution of the
recurrence~\eqref{eq:oe-recurrence} in $n$, with $y_r=0$ and, by
Lemma~\ref{lem:oe-casoratian},
$y_{r+1}=b_rc_{r+1}-c_rb_{r+1}=-W_{r+1}=1/(r+1)^3$.
Let $\hat y$ be the normalized restart solution, $\hat y_r=0$,
$\hat y_{r+1}=1$, so $y_n=(r+1)^{-3}\hat y_n$ by linearity.  We claim
\begin{equation}\label{eq:coinc-normalized-restart}
 \hat y_{r+h}=\frac{N_h(r)}{\prod_{j=2}^{h}(r+j)^3}\qquad(r+h\le M),
\end{equation}
which
gives~\eqref{eq:coinc-gap-identity} after multiplying by $(r+1)^{-3}$.
For $h=1$, both sides are $1=N_1$.  For $h=2$, the recurrence at $n=r+1$
reads $(r+2)^3\hat y_{r+2}=P(r+1)\hat y_{r+1}-(r+1)^3\hat y_r=P(r+1)=N_2(r)$.
Assume~\eqref{eq:coinc-normalized-restart} for $h-1$ and $h$.  The
recurrence at $n=r+h$, multiplied by $\prod_{j=2}^{h}(r+j)^3$, reads
\[
 (r+h+1)^3\!\!\prod_{j=2}^{h}(r+j)^3\,\hat y_{r+h+1}
 =P(r+h)N_h(r)-(r+h)^6N_{h-1}(r)
 =N_{h+1}(r),
\]
by the gap recurrence of Lemma~\ref{lem:gap-poly}; this
is~\eqref{eq:coinc-normalized-restart} for $h+1$.  The zero
characterization is the dictionary~\eqref{eq:oe-dictionary} together with
the invertibility of the product.
\end{proof}

For $1\le H\le M-1$ define the index set and the value counts
\[
 \mathcal S_H=\{(r,h):1\le h\le H,\ 1\le r\le M-h\},
 \qquad
 n_H(a)=\#\{(r,h)\in\mathcal S_H:\Delta_{r,h}=a\},
\]
with $S:=\#\mathcal S_H=\sum_{h\le H}(M-h)$, and the coincidence second
moment
\[
 \mathcal N(H)=\sum_{a\in\mathbb F_p}n_H(a)^2
 =\#\{\bigl((r,h),(r',h')\bigr)\in\mathcal S_H^2:
 \Delta_{r,h}=\Delta_{r',h'}\}.
\]

\begin{conjecture}[coincidence target, $\mathrm{COINC}_\beta$]
\label{conj:coinc-beta}
Let $\tfrac12<\beta\le1$.  There is an absolute constant $K$ such that,
for every prime $p\ge7$, with $H=\lceil p^\beta\rceil\wedge(M-1)$,
\begin{equation}\label{eq:coinc-hypothesis}
 \mathcal N(H)\le\frac{S^2}{p}+K\,S .
\end{equation}
This is open for every $\beta>\tfrac12$; it asserts that the values
$\Delta_{r,h}$ equidistribute in second moment at scale $H$.
\end{conjecture}

\begin{lemma}[centered-count inequality]\label{lem:coinc-centered}
Unconditionally, for every $1\le H\le M-1$,
\begin{equation}\label{eq:coinc-centered}
 \sum_{h=1}^{H}C_h\;\le\;\frac Sp+\sqrt{\mathcal N(H)-\frac{S^2}p\,}.
\end{equation}
\end{lemma}

\begin{proof}
By Lemma~\ref{lem:coinc-gap-identity}, $\sum_{h\le H}C_h=n_H(0)$.  Since
$\sum_an_H(a)=S$,
\[
 \Bigl(n_H(0)-\frac Sp\Bigr)^2
 \le\sum_{a\in\mathbb F_p}\Bigl(n_H(a)-\frac Sp\Bigr)^2
 =\mathcal N(H)-\frac{S^2}p .
\]
\end{proof}

\begin{theorem}[conditional energy exponents from coincidence counting]
\label{thm:coinc-energy}
Assume Conjecture~\ref{conj:coinc-beta} for one value
$\tfrac12<\beta\le1$ with constant $K$.  Then for $\tfrac12<\beta<1$,
\[
 \mathcal F_p\le(1+o(1))\,p^{2-\beta},
\]
and for $\beta=1$,
\[
 \mathcal F_p\le(3+2\sqrt K)\,p ,
\]
i.e.\ E1 holds.  In particular any single $\beta>\tfrac12$ strictly
improves the unconditional exponent $3/2$ of
Theorem~\ref{thm:oe-energy-three-halves}.
\end{theorem}

\begin{proof}
Let $H=\lceil p^\beta\rceil\wedge(M-1)$.
By Lemma~\ref{lem:coinc-centered} and~\eqref{eq:coinc-hypothesis},
using $S\le MH$,
\begin{equation}\label{eq:coinc-linear-sum}
 \sum_{h\le H}C_h\le\frac{MH}p+\sqrt{K\,MH}\le H+\sqrt{K\,pH}.
\end{equation}
Partition $I_p$ into $\lceil M/(H+1)\rceil$ consecutive blocks of length
at most $H+1$; the proof of Theorem~\ref{thm:oe-energy-three-halves},
verbatim with block length $H+1$ in place of $H$, gives
\[
 \mathcal F_p\le
 \Bigl\lceil\frac M{H+1}\Bigr\rceil
 \Bigl(M+2\sum_{h=1}^{H}C_h\Bigr)
 \le\Bigl(\frac MH+1\Bigr)
 \Bigl(M+2H+2\sqrt{KpH}\Bigr).
\]
For $\tfrac12<\beta<1$ the right side is
$p^{2-\beta}\bigl(1+2\sqrt K\,p^{(\beta-1)/2}+O(p^{\beta-1}+p^{-\beta}+p^{2\beta-2})\bigr)
=(1+o(1))p^{2-\beta}$, since every correction exponent is negative for
$\tfrac12<\beta<1$.  For $\beta=1$ one block suffices ($H=M-1$):
$\mathcal F_p\le M+2\sum_{h\le M-1}C_h\le M+2(M-1)+2\sqrt{K\,p\,(M-1)}
\le(3+2\sqrt K)p$.
\end{proof}

\begin{proposition}[position in the target hierarchy]
\label{prop:coinc-position}
Fix a scale $1\le H\le p^{1-\delta}$, $\delta>0$.  Then, with $B_t(H)$ as
in Conjecture~\ref{conj:oe-2D-sqrt} extended to the index set
$\mathcal S_H$,
\begin{equation}\label{eq:coinc-parseval}
 \sum_{t\ne0}\bigl|B_t(H)\bigr|^2
 =p\,\mathcal N(H)-S^2 .
\end{equation}
Consequently the square-root cancellation
$|B_t(H)|\le C(pH)^{1/2}p^{\varepsilon}$ for \emph{every} $t\ne0$
(GAP--2D--SQRT at scale $H$) implies~\eqref{eq:coinc-hypothesis} with
$K=C^2p^{2\varepsilon}(1+o(1))$: the coincidence target is precisely the
$L^2$-average over $t$ of GAP--2D--SQRT, and is therefore the weakest
member of the sufficient family at its scale.  It is not comparable to
ATR or GPRV in either direction by any argument recorded in this paper.
Two honesty clauses.  First, the exact partition
$\mathcal N(H)=S+A(A-1)+R_{\ne0}$, with $A=\sum_{h\le H}C_h$ and
$R_{\ne0}$ the distinct-index coincidences at nonzero values, shows that
the centered variance \emph{contains} the term $(A-S/p)^2$ as one
nonnegative coordinate: Conjecture~\ref{conj:coinc-beta} is a
reformulation with a scale dial and a weakest-sufficient position, not a
reduction in difficulty --- any proof of it must control the zero fiber
at exactly the scale of~\eqref{eq:coinc-linear-sum}.  Second, E1 does not
imply $\mathrm{COINC}_\beta$ and $\mathrm{COINC}_\beta$ at a truncated
scale does not imply E1; the two are genuinely different statements about
the same value distribution.
\end{proposition}

\begin{proof}
Orthogonality of additive characters gives
$\sum_{t\in\mathbb F_p}|B_t|^2=p\sum_an_H(a)^2=p\,\mathcal N(H)$, and the
$t=0$ term is $S^2$.  For the implication, sum the pointwise bound over
$t\ne0$: $p\,\mathcal N(H)-S^2\le(p-1)\,C^2pH\,p^{2\varepsilon}$, and
$S=\sum_{h\le H}(M-h)\ge H(M-H)=(1-o(1))\,pH$ for $H\le p^{1-\delta}$,
so $\mathcal N(H)-S^2/p\le C^2p^{2\varepsilon}(1+o(1))\,S$.
\end{proof}

\begin{remark}[why $\beta>\tfrac12$ is the battle line]
\label{rem:coinc-battle-line}
For $H\le p^{1/2}$ the bound~\eqref{eq:coinc-linear-sum} gives
$\sum_{h\le H}C_h\ll(pH)^{1/2}$, which through the block argument
reproduces the exponent $3/2$ and nothing better; this is the exact
content of the crossover in Theorem~\ref{thm:coinc-energy}.  Moreover
$(pH)^{1/2}\ge H^2$ exactly when $H\le p^{1/3}$, so
below $p^{1/3}$ the conclusion of~\eqref{eq:coinc-linear-sum} is weaker
than the trivial degree bound $\sum_{h\le H}C_h\le\tfrac32H^2$: partial
progress on Conjecture~\ref{conj:coinc-beta} has content only for
$\beta>\tfrac13$, and improves the energy record only for
$\beta>\tfrac12$.  Pairwise curve counting cannot reach even the first
threshold: the difference correspondence
$\{\Delta_{r,h}=\Delta_{r',h'}\}$ has bidegree $(3h,3h')$, hence generic
arithmetic genus $O(hh')$ and Weil error $O(hh'\sqrt p)$ per pair, so the
summed error $O(H^4\sqrt p)$ fits the $O(pH)$ budget only for
$H\le p^{1/6}$ (or $H\le p^{1/4}$ if the continuant structure forces a
genus collapse to $O(h{+}h')$, which is not known).  On the other side,
for $H>p^{1/2}$ the degrees $3(h-1)$ exceed $\sqrt p$ and pointwise Weil
bounds carry errors at least as large as the main terms.  Every route
through fixed-pair algebraic geometry is therefore dead in the regime
where Conjecture~\ref{conj:coinc-beta} has content.  This is one more
independent arrival at the family-compatibility frontier of
Remark~\ref{rem:oe-chebotarev-limitation}.
\end{remark}

\begin{remark}[the exponent $3/2$ is the set-genericity threshold]
\label{rem:coinc-vinh}
The pairs $(n,m)$ with $\pi(n)=\pi(m)$ are exactly the incidences between
the point set $\{v_n\}$ and the line set $\{\mathrm{span}(v_m)\}$ in
$\mathbb F_p^2$, and Vinh's finite-field incidence theorem gives, for
\emph{any} $M$ points and $M$ lines, at most $M^2/p+p^{1/2}M$ incidences
--- with $M\sim p$ this is again $O(p^{3/2})$.  Thus
Theorem~\ref{thm:oe-energy-three-halves} attains precisely the bound that
holds for a set-generic configuration, and \emph{no} argument that uses
the orbit only as a set of directions can pass $3/2$: E1 and every
$\beta>\tfrac12$ instance of Conjecture~\ref{conj:coinc-beta} are
assertions about the arithmetic ordering of the sequence
$n\mapsto v_n$, invisible to incidence geometry.
\end{remark}

\begin{remark}[numerical status only]\label{rem:coinc-empirical}
For $p\in\{499,997,1999,4001,7919\}$ and $H=\lceil p^{0.66}\rceil$, the
measured value of $\bigl(\mathcal N(H)-S^2/p\bigr)/S$ is
$1.24,\,1.08,\,1.08,\,1.08,\,1.05$: bounded and slowly decreasing in $p$,
consistent with Conjecture~\ref{conj:coinc-beta} at $\beta=2/3$ with
$K$ close to $1$.  Identity~\eqref{eq:coinc-gap-identity} was verified
exactly at $300$ random pairs $(r,h)$ for each of the five primes.  These
are numerical data, not theorems, and none is used above.
\end{remark}
